%% ma: L12-B08-0116
%% lop: 12
%% chuong: 2
%% bai: 8
%% dang: 12.8.26
%% loai: TN4
%% muc: TH
%% dap_an: "C"
%% nguon: "CĐ2.3.pdf giáo viên gửi 27/10/2026 (M07, Lê Minh Khôi – Nguyễn Anh Khôi, Chương II Chủ đề 3 Biểu thức tọa độ của các phép toán vectơ), Câu 28"
%% kien_thuc: "Bài 8 (khoảng cách hai điểm)"
%% trang_thai: chinh-thuc
\begin{ex}
Trong không gian $Oxyz$, tọa độ của điểm $M$ thuộc $Ox$ và cách đều hai điểm $A(4;2;-1)$ và $B(2;1;0)$ là
\choice{$M(-4;0;0)$}{$M(5;0;0)$}{\True $M(4;0;0)$}{$M(-5;0;0)$}
\loigiai{$M(x;0;0)$: $(x-4)^2+5=(x-2)^2+1\Leftrightarrow-4x=-16\Leftrightarrow x=4$.}
\end{ex}
